% Direct-authority diplomatic French for printed p.185.
\clearpage
\SourcePage{185}

\LayerTarget{tohoku:chapter:IV}
\begin{center}
{\large\bfseries Chapitre IV\quad Les Ext de faisceaux de modules.}
\end{center}

\LayerTarget{tohoku:section:4.1}
\SectionHeading{4.1. Les foncteurs $\mathrm{Hom}_O(A,B)$ et
$\mathbf{Hom}_O(A,B)$.}

\LayerTarget{tohoku:definition:hom-o-and-local-hom}
Soit $X$ un espace topologique muni d'un faisceau $O$ d'anneaux avec unité,
et soit $C^O$ la catégorie abélienne des $O$-Modules (à gauche) sur $X$
(cf. 3.1). Si $A$ et $B$ sont deux $O$-modules, nous désignons par
$\mathrm{Hom}_O(A,B)$ le groupe des $O$-homomorphismes du premier dans le
second (ce groupe est en fait un module sur le centre de $\Gamma(O)$). Si $U$
est une partie quelconque de $X$, on posera
\[
\mathrm{Hom}_O(U;A,B)=\mathrm{Hom}_{O|U}(A|U,B|U)
\]
(où $F|U$ désigne, comme d'habitude, la restriction d'un faisceau $F$ à $U$).

\LayerTarget{tohoku:definition:sheaf-hom-o}
Bien entendu, si $V\subset U$, on a un homomorphisme naturel
$\mathrm{Hom}_O(U;A,B)\longrightarrow\mathrm{Hom}_O(V,A,B)$, et on vérifie
immédiatement que si on se restreint aux parties ouvertes $U$ de $X$, les
groupes $\mathrm{Hom}_O(U;A,B)$ forment un \textit{faisceau} sur $X$, noté
$\mathbf{Hom}_O(A,B)$. Ainsi par définition, on a

\LayerTarget{tohoku:eq:4.1.1}
\[
\tag{4.1.1}
\Gamma(U,\mathbf{Hom}_O(A,B))=\mathrm{Hom}_O(U;A,B)
\]
et en particulier, faisant $U=X$ :

\LayerTarget{tohoku:eq:4.1.2}
\[
\tag{4.1.2}
\mathrm{Hom}_O(A,B)=\Gamma(\mathbf{Hom}_O(A,B)).
\]

\LayerTarget{tohoku:remark:sheaf-hom-o:center-module}
D'ailleurs, $\mathbf{Hom}_O(A,B)$ peut aussi être considéré comme un faisceau
de modules sur le centre $O^\natural$ de $O$.

\LayerTarget{tohoku:remark:hom-o:left-exactness}
Rappelons que $\mathrm{Hom}_O(A,B)$ est un \textit{foncteur additif exact à
gauche} en les deux arguments $A,B\in C^O$, à valeurs dans les groupes
abéliens. On en conclut que $\mathbf{Hom}_O(A,B)$ peut aussi être regardé comme
un foncteur exact à gauche en les deux arguments $A,B\in C^O$, à valeurs dans
la catégorie $C^X$ des faisceaux abéliens sur $X$, (ou même à valeurs dans la
catégorie $C^{O^\natural}$). Comme d'habitude, tous les homomorphismes que nous
écrirons seront ``fonctoriels''. Ce numéro donne quelques propriétés
auxiliaires des foncteurs envisagés, préliminaires à l'étude de leurs foncteurs
dérivés.

Soient $A,B\in C^O$, et soit $x\in X$. Pour tout ouvert $U$ contenant $x$, on
a un homomorphisme évident
$\mathrm{Hom}_O(U;A,B)\longrightarrow\mathrm{Hom}_{O(x)}(A(x),B(x))$, d'où par
passage à la limite inductive sur les voisinages ouverts de $x$, un
homomorphisme naturel

\LayerTarget{tohoku:eq:4.1.3}
\[
\tag{4.1.3}
\mathbf{Hom}_O(A,B)(x)\longrightarrow\mathrm{Hom}_{O(x)}(A(x),B(x))
\]
que nous allons étudier.

\LayerTarget{tohoku:definition:finite-type-at-x}
Nous dirons que $A$ est \textit{de type fini en $x$} si on peut trouver un
voisinage ouvert $U$ de $x$ tel que $A|U$ soit isomorphe à un quotient de
$(O|U)^n$ (où $n$ est un entier fini $>0$) ;

\LayerTarget{tohoku:definition:pseudo-coherent-at-x}
$A$ est dit \textit{pseudo-cohérent en $x$} si on peut trouver un voisinage
ouvert $U$ de $x$, et \textit{sur $U$} une suite exacte d'homomorphismes
$O^m\longrightarrow O^n\longrightarrow A\longrightarrow0$ ($m,n$ entiers
finis $>0$). Nous dirons que $A$ est de type fini (resp. pseudo-cohérent) s'il
l'est en tout point. Ceci posé :

\LayerTarget{tohoku:proposition:4.1.1}
\noindent\textsc{Proposition 4.1.1.} \textit{L'homomorphisme (4.1.3) est un
monomorphisme si $A$ est de type fini en $x$, et un isomorphisme si $A$ est
pseudo-cohérent en $x$.}

\LayerTarget{tohoku:proof:proposition-4.1.1:p1}
Supposons $A$ de type fini en $x$, alors, en se restreignant au besoin à un
voisinage convenable de $x$, on a une suite exacte
$O^n\longrightarrow A\longrightarrow0$ d'où une suite

\TohokuContinuation{186}
